\section{Appendix: API Reference}

\subsection{Core Classes}
\begin{itemize}
    \item \texttt{HardwareBackend} (abstract): Defines the interface for all backends.
    \item \texttt{HardwareFactory}: Singleton factory for backend instantiation and selection.
    \item \texttt{BackendRegistry}: Maintains the prioritized list of available backends.
    \item \texttt{HardwareConfig}: Root configuration dataclass with sub-configs for each backend.
\end{itemize}

\subsection{Key Methods of HardwareBackend}
\begin{itemize}
    \item \texttt{initialize(config: Dict) -> bool}: Initialize the backend.
    \item \texttt{create\_continuous\_field(dim: int) -> np.ndarray}: Create a normalized field.
    \item \texttt{field\_update(field, dt, past, present, future) -> Dict}: Evolve the field.
    \item \texttt{compute\_interference(a, b) -> float}: Compute overlap between two fields.
    \item \texttt{find\_stable\_patterns(fields, threshold) -> List[Dict]}: Find mutually coherent field pairs.
    \item \texttt{measure\_coherence(fields) -> float}: Global coherence metric.
    \item \texttt{cleanup()}: Release backend resources.
\end{itemize}

\subsection{Adding a New Backend}
To add a new backend (e.g., for a neuromorphic chip):
\begin{enumerate}
    \item Create a new class inheriting from \texttt{HardwareBackend}.
    \item Implement all abstract methods.
    \item Register the backend in \texttt{BackendRegistry} with an appropriate priority and required packages.
    \item Optionally, add a configuration dataclass in \texttt{config.py}.
\end{enumerate}
The factory will automatically detect and initialize the new backend when its dependencies are satisfied.